\documentclass[conf]{new-aiaa}
\usepackage[utf8]{inputenc}
\usepackage{graphicx}
\usepackage{amsmath}
\usepackage{siunitx}
\usepackage{longtable,tabularx}
\usepackage{booktabs}
\usepackage{float}
\usepackage[section]{placeins}
\usepackage{xcolor}
\usepackage{array}
\usepackage{xurl}
\newcommand{\gpc}[1]{\textcolor{teal}{\small\textsf{[#1]}}} % GP-class tag
\graphicspath{{figures/}}
\emergencystretch=2.5em
\AtBeginDocument{%
  \setlength{\intextsep}{5pt plus 2pt minus 1pt}%
  \setlength{\textfloatsep}{6pt plus 2pt minus 2pt}%
  \setlength{\abovecaptionskip}{3pt}%
  \setlength{\belowcaptionskip}{1pt}%
}

\title{Layer-A Geometric-Programming Sizing of a Single Oblique-Wing Trans-Medium Vehicle: CFD-Calibrated Surrogates and a Buoyancy-Limited Payload Frontier}
\author{Yiyang I. Zhang\footnote{Independent Researcher, Winchester College.}}
\affil{Winchester College, Winchester, United Kingdom}

\begin{document}
\maketitle

\section{Methodology}

\subsection{Modelling framework and the scope of Layer~A}
The vehicle is a slender, missile-shaped hull carrying a single straight wing on one central
vertical pivot. In air the wing is deployed perpendicular to the body to generate lift; before the
vehicle enters the water the wing rotates flat against the hull, so that underwater the vehicle is a
clean body of revolution whose small residual overhang provides a hydrodynamic trim force. Sizing
such a vehicle couples three normally separate disciplines---aerodynamics, hydrodynamics, and
structures---and two operating media whose dynamic pressures differ by more than two orders of
magnitude. To keep the coupled problem tractable and, crucially, to guarantee that the optimum found
is the \emph{global} one, the design is decomposed into three layers. Layer~A, the subject of this
paper, is a geometric program (GP) that fixes every continuous sizing variable for a prescribed
wing-skew angle and a prescribed buoyancy state. Layer~B, a Bayesian optimiser that schedules the
wing-skew trajectory during the dive~\cite{long2025entire}, sits above Layer~A and is not treated here. A computational
fluid dynamics (CFD) campaign underpins both layers by supplying the aerodynamic and hydrodynamic
coefficients that cannot be written analytically.

A geometric program is an optimisation problem in which a posynomial objective is minimised subject
to posynomial inequality constraints and monomial equality constraints. A \emph{monomial} is a
product of powers of the design variables with a positive leading coefficient; a \emph{posynomial}
is a sum of monomials. The defining property that makes the formulation attractive is that, under the
change of variables $u_i\mapsto\ln u_i$, every GP becomes convex, so a solver returns the global
optimum---not merely a local one~\cite{boyd2007tutorial}---together with the sensitivity of the objective to every
constraint. The engineering price of this guarantee is that each governing relation must be cast into
posynomial or monomial form; the methodology below is largely the discipline of paying that price
honestly.

\subsection{CFD design of experiments}\label{sec:cfd}
The coefficients that Layer~A cannot derive from first principles are the aerial and underwater drag
polars, the lift-curve slope and its variation with wing skew, the stowed-wing trim lift, and the
cavitation suction peak. These, and only these, define the CFD design of experiments. Two campaigns
were run in OpenFOAM~v2606 with the steady incompressible solver \texttt{simpleFoam} and the
$k$--$\omega$ SST turbulence model. The aerial campaign meshed all ten wing positions from stowed
($\Lambda=90^\circ$) to deployed ($\Lambda=0^\circ$) at $10^\circ$ intervals inside a free-air domain
of $17\times12\times12$~m, and swept the angle of attack from $-4^\circ$ to $+14^\circ$ by rotating
the inlet velocity vector rather than re-meshing, which saved roughly nine-tenths of the meshing
effort. The underwater campaign reused the same meshes with water properties, added a speed sweep of
$2$ to $20$~m/s at the deployed and stowed extremes, and treated operating depth as a
post-processing parameter, since in an incompressible solution the flow field is independent of the
absolute pressure level and depth enters only through the cavitation number. Every force coefficient
was referenced consistently to the wing planform area, and convergence was judged by coefficient
stability---the mean of the final fifty iterations was required to differ from the preceding fifty by
less than one percent---rather than by a fixed residual threshold. As an independent check on the
whole pipeline, the stowed induced-drag factor obtained in air and in water, at Reynolds numbers that
differ by a factor of about $1.5$, agreed to within $7\%$; because the drag coefficient is
dimensionless and, at matched Reynolds number, depends only on geometry, this agreement confirms that
the two campaigns are internally consistent.

\subsection{GP-compatible surrogate fitting}
Each CFD relationship was reduced to a GP-legal surrogate with the \texttt{gpfit} workflow of Hoburg~et~al.~\cite{hoburg2016data}. Where a
relationship is affine on logarithmic axes it was fitted as a single monomial by ordinary least
squares in log space, which is mathematically identical to a one-term maximum-affine fit and is
numerically robust; where a relationship carries genuine curvature, a multi-term softmax-affine
posynomial was fitted instead. The lift-curve slope varies with wing skew as
\begin{equation}
C_{L\alpha}(\Lambda) = 0.847\,(\cos\Lambda)^{1.105}\ \text{rad}^{-1},
\qquad(\text{rms}_{\log}=0.011),
\label{eq:clalpha}
\end{equation}
whose exponent of $1.105$ is within eleven percent of the value of unity predicted by the
sweep-independence principle, so the fit does more than interpolate: it confirms the analytic
expectation that only the velocity component normal to the wing generates lift. The induced-drag
factor rises steeply as the wing folds and its effective aspect ratio collapses,
\begin{equation}
k(\Lambda) = 0.631\,(\cos\Lambda)^{-1.352}\ ,
\qquad(\text{rms}_{\log}=0.013),
\label{eq:kschedule}
\end{equation}
climbing from $0.64$ at full deployment to $22.9$ when stowed---a thirty-six-fold increase that is
the direct signature of the aspect-ratio collapse in the relation $k=1/(\pi e A)$. The underwater
zero-lift drag scales with Reynolds number as $C_{D0,w}=0.185\,\mathrm{Re}^{-0.112}$, an exponent
that sits between the $-0.2$ of pure turbulent skin friction and the zero of pure form drag,
consistent with a body whose drag is a mixture of the two. The cavitation suction peak was found to
be essentially constant, as a dimensionless quantity must be: for the stowed configuration, whose
trim lift keeps it near zero lift coefficient, $-C_{p,\min}\approx1.16$, and this constant is used as
the conservative cavitation gate.

\section{The geometric-programming model}

\subsection{Objective and requirements}
The objective is to \emph{maximize payload} for the required ranges:
\begin{equation}
\text{minimize}\quad m_{\text{pay}}^{-1},\qquad R_a\ge R_a^{\min},\ \ R_w\ge R_w^{\min}.
\gpc{monomial objective}
\end{equation}

\subsection{Aerial cruise: steady level flight}
With the wing deployed, lift balances weight, thrust balances drag~\cite{hoburg2014geometric}:
\begin{align}
m g &\le \tfrac12\rho_a V_a^2 S\,C_{L,a}, \gpc{monomial} \\
D_a &\ge \tfrac12\rho_a V_a^2 S\,C_{D,a}, \gpc{monomial (relaxed)} \\
P_a &\ge D_a V_a / \eta_a. \gpc{monomial}
\end{align}
These are the aircraft steady-flight relations of Ref.~\cite{hoburg2014geometric}, Eqs.~(15)--(18).

\subsection{Aerodynamic drag build-up}
The drag coefficient decomposes into induced, profile, and form
contributions~\cite{hoburg2014geometric,hoerner1965fluid}:
\begin{align}
C_{D,a} &\ge \frac{C_{L,a}^2}{\pi e A}
           + C_{Dp}(C_{L,a},\mathrm{Re},\Lambda)
           + C_{D0,a}, \gpc{posynomial}\\[4pt]
C_{D0,a} &= C_f\,(1+k)\,\frac{S_{\text{wet}}}{S},
           \qquad C_f = \frac{0.074}{\mathrm{Re}^{0.2}},
           \qquad \mathrm{Re} = \frac{\rho_a V_a c}{\mu}, \gpc{monomial}\\[4pt]
(1+k) &= 1 + 1.5\left(\frac{d}{L}\right)^{3/2}
        +     7.0\left(\frac{d}{L}\right)^{3}. \gpc{posynomial}\label{eq:formfactor}
\end{align}
Equation~\eqref{eq:formfactor} is Hoerner's~\cite{hoerner1965fluid} Eq.~(28) for streamlined bodies of
revolution. The term $1.5(d/L)^{3/2}$ captures the supervelocity increment (Eq.~26, $\Delta q_{\text{av}}/q
\propto(d/L)^{3/2}$); the $7(d/L)^3$ term is the separation/pressure component, negligible for
$d/L\lesssim0.2$. For our hull ($d/L\approx0.10$): $(1+k)\approx1.054$, i.e.\ a
\emph{5.4\%} viscous drag adder from thickness — a far smaller contribution than the $\sim$20\%
assumed in the design brief, which conservatively included appendages. The SLA (service-life-allowance)
roughness/appendage increment is added as an empirical factor $f_{\text{SLA}}$, i.e.\
$(1+k)_{\text{eff}}=f_{\text{SLA}}(1+k)_{\text{naked}}$. The turbulent flat-plate law
$C_f=0.074\,\mathrm{Re}^{-0.2}$~\cite{hoburg2014geometric} is monomial and used in place of the ITTC
1957 log-law~\cite{myring1976theoretical,allen2000propulsion}; for $\mathrm{Re}\approx
2\times10^{7}$--$3\times10^{6}$ (water/air on $L\approx 1$~m) the power-law fit error is $<2\%$.\par
The profile-drag term $C_{Dp}(C_{L,a},\mathrm{Re},\Lambda)$ is \emph{not analytic} and must be supplied
as a fitted posynomial from CFD (Sec.~\ref{sec:cfd}), following the gpfit workflow of
Ref.~\cite{hoburg2016data}. The $\Lambda$-dependence of the Oswald factor $e(\Lambda)$ is also CFD-fitted;
at $\Lambda=0$ the analytic prior is $e\approx0.9$ from lifting-line theory~\cite{hoerner1985fluid}.

\subsection{Propulsion and endurance (electric)}
For a battery-electric vehicle the fuel-burn/Breguet model of Ref.~\cite{hoburg2014geometric} is
replaced by a constant-mass energy balance; endurance $t$ equals stored energy over
power~\cite{allen2000propulsion,rutherford2008southampton,lin2026energies}:
\begin{align}
E_b   &= e_b\, m_b,                  & R_a   &= V_a t_a,\quad R_w = V_w t_w, \gpc{monomial}\\
E_b   &\ge P_a t_a + P_w t_w + P_{\text{hotel}}t_{\text{total}}, \gpc{posynomial}\\
\eta_a&\le \eta_{\text{motor}}\,\eta_{\text{prop}}. \gpc{monomial}
\end{align}
For the aerial leg, $\eta_{\text{prop}}$ follows the actuator-disk model of
Ref.~\cite{hoburg2014geometric}, Eq.~(37), with $\eta_i\le 2/(1+\sqrt{1+T/(\frac12\rho V^2
A_{\text{prop}})})$. For the underwater leg, the optimal propeller efficiency is approximately
constant $\eta_{\text{prop,uw}}\approx0.81$~\cite{allen2000propulsion}. System-level Li-ion
specific energy is $\approx100$--$180$~Wh/kg (cell: $190$--$300$~Wh/kg)~\cite{lin2026energies};
REMUS field data~\cite{allen2000propulsion} confirm $1000$~Wh over $20$~h at $1.5$~m/s ($P
\approx50$~W). Hotel load $\approx 200$--$220$~W (survey-class AUV)~\cite{lin2026energies}; for
this one-off sprint design, factor it into the fixed-mass term or drop if payload power is
dominant.

\subsection{Mass build-up}
\begin{align}
m &\ge m_{\text{pay}} + m_s + m_b + m_{\text{prop}} + m_{\text{fixed}}, \gpc{posynomial}\\
m_{\text{prop}} &\le c_m\,P_{\max}. \gpc{monomial (power-law)}\\
m_s &\ge m_{\text{skin}} + m_{\text{fins}} + m_{\text{pivot}}. \gpc{posynomial}
\end{align}
If the hull is unstiffened (plausible at $100$~mm diameter), $m_{\text{skin}} = \rho_{\text{mat}}\,\pi
d\,L\,t_s$ with $t_s$ sized by the buckling gate (Sec.~\ref{sec:structure}). The power-law mass
coefficient $c_m$ (kg/kW) for a brushless DC motor + controller is $\approx 0.1$--$0.3$~kg/kW
(representative UAV/UUV values); verify against specific vendor data.

\subsection{Underwater cruise: buoyancy and wing trim}\label{sec:buoyancy}
Underwater the wing is stowed. Buoyancy supports most of the weight; the stowed overhang provides a
hydrodynamic trim force that closes the vertical balance. With $B=\rho_w g\nabla$ and the imbalance
fraction $\xi=(W-B)/B$ \emph{swept as a parameter}, the balance is monomial per solve:
\begin{align}
B   &= \rho_w g\, \nabla, \gpc{monomial}\\
m g &= B\,(1+\xi),        \gpc{monomial in $\xi$ (fixed per solve); SP if $\xi$ free}\\
L_{\text{tr}} &= \xi B = \tfrac12\rho_w V_w^2 S_{\text{stow}}\,C_{L,\text{uw}}. \gpc{monomial}
\end{align}
Sweeping $\xi$ from slightly negative (positively buoyant, wing trims down) to positive (heavy,
wing lifts — payload exceeds the buoyancy ceiling) traces a payload--trim-drag frontier. The trim
lift $L_{\text{tr}}$ induces drag $D_{\text{tr}}\ge k_{\text{uw}} L_{\text{tr}}^2/
(\tfrac12\rho_w V_w^2 S_{\text{stow}})$ that feeds the underwater power term.
The stowed wing's hydrodynamic behaviour is self-compensating: as the wing folds flat its aspect
ratio collapses $A\sim3.3 \to 0.1$ (Jones slender-body limit $\sim\pi A/2$~\cite{hoerner1985fluid}),
which reduces $C_{L\alpha}$ by roughly the same factor ($\sim22$) as the water/air $q$-ratio
($\sim130$) divided by the area reduction ($\sim3$): $130/(3\cdot22)\approx1.9$. This
near-cancellation — $\textit{not}$ an exact equality — is the analytic backbone of the
self-compensation thesis in the design brief and is tested directly by the HYDRO CFD campaign.
The added-mass correction for vertical motion is negligible in the axial direction
($k_a\approx0.03$ at $L/d\approx10$~\cite{newman2018marine}) but may be $\sim$95\% of the
displaced mass laterally~\cite{newman2018marine}; on entry the effective mass $(m+m_a)$ replaces
$m$ in the impact constraint (Sec.~\ref{sec:entry}) and in any lateral-belly-flop transient (not
modelled here).

\subsection{Underwater hydrodynamic drag}
For the slender body of revolution,
\begin{align}
D_w      &= \tfrac12\rho_w V_w^2 S_{\text{wet}}\,C_{D,w}, \gpc{monomial}\\
C_{D,w}  &\ge C_f\,(1+k) + k_{\text{uw}}\,
            \frac{L_{\text{tr}}^2}{(\tfrac12\rho_w V_w^2)^2 S_{\text{stow}}^2}, \gpc{posynomial}
\end{align}
with $(1+k)$ from Eq.~\eqref{eq:formfactor} (same form factor, water Reynolds number). The trim-lift
induced-drag coefficient $k_{\text{uw}}$ comes from CFD (stowed wing, $\Lambda=90^\circ$). At
$d/L=0.10$, the minimum-drag thickness ratio $\approx 0.18$--$0.20$~\cite{hoerner1965fluid}; our
hull falls on the thinner (friction-dominated) side.

\subsection{Cavitation margin}\label{sec:cavitation}
The cavitation inception condition~\cite{brennen2014cavitation} is
\begin{equation}
-C_{p,\min}\ \le\ \sigma_i,
\qquad \sigma = \frac{p_{\text{atm}} + \rho_w g h - p_v}{\tfrac12\rho_w V_w^2},
\qquad p_v(20^\circ\text{C})\approx 2.34\ \text{kPa}.
\gpc{posynomial $\le$ monomial}
\end{equation}
At the surface ($h=0$, $V_w=20$~m/s): $\sigma\approx0.495$, consistent with
the design brief. Amromin and Rozhdestvensky~\cite{amromin2022cavitation} show that at hull Reynolds
numbers $\mathrm{Re}\gtrsim10^7$ (our $\mathrm{Re}\approx2\times10^7$) the difference
$|\min C_p|-\sigma_i\to0$, i.e.\ the $\sigma_i=-C_{p,\min}$ rule is increasingly accurate. The
constraint is therefore conservative. Operating below this bound keeps the vehicle clear of the supercavitating regime addressed for dedicated high-speed underwater bodies~\cite{dzielski2003benchmark,zou2023optimized}, which lies outside the present cavitation-limited design; the suction peak $-C_{p,\min}$ is the \emph{CFD deliverable}
(fitted vs $C_L$ or $\alpha$ from the HYDRO campaign). The GP enters it as
$\tfrac12\rho_w V_w^2(-C_{p,\min})+p_v\le p_{\text{atm}}+\rho_w g h$, which is posynomial in $V_w$
(with $-C_{p,\min}$ pre-fitted as a monomial/posynomial in $C_L$).

\subsection{Pressure-hull structure and buckling}\label{sec:structure}
For operating depth $h$, the thin-shell hoop stress (unstiffened tube) gives a monomial gauge
constraint~\cite{shinoka2024structural}:
\begin{equation}
\sigma_h = \frac{p_d\, d}{2 t_s} \le \sigma_y \ \Rightarrow\
p_d \le \frac{2\sigma_y t_s}{d},
\qquad p_d = \rho_w g h. \gpc{monomial}
\end{equation}
External-pressure buckling adds a second limit. The Bryant overall-elastic-buckling collapse
pressure~\cite{shinoka2024structural} is
\begin{equation}
p_{\text{cr}} = \frac{E\,t_s/a}{\,n^2-1+\lambda^2/2\ }
                \left[\frac{1}{(n^2+ \lambda^2)^2}\right] + \cdots
                \gpc{signomial (discrete $n$) / non-GP}
\end{equation}
where $\lambda=\pi a/L$, $a=d/2$ the radius, $n$ the integer buckling lobe number ($\ge2$).
The full Bryant form (not shown in full here) contains an additional $\lambda^4$ factor in the
numerator; this displayed sketch omits it — the equation is illustrative and will be replaced by a
fitted monomial in the GP implementation.
This is \emph{signomial} (posynomial denominator inverted) and has a discrete inner minimization over
$n$. For a thin unstiffened tube at our small diameter ($100$~mm), the recommended GP path is to
\emph{factor out} the $n$-minimization into an offline evaluation and fit a monomial
$p_{\text{cr}} \approx C\,E\,(t_s/a)^\alpha\,(a/L)^\beta$ over the $(t_s/a,\ L/a)$ range,
following the Windenburg--Trilling closed form. The operational safety factor is
$\mathrm{SF}_{\text{overall}}\ge2.0$~\cite{shinoka2024structural}; the constraint in the GP is
then $p_d \le p_{\text{cr}}/\mathrm{SF}$, which remains monomial after the fit.
Material: for a first pass use aluminium 6061-T6 ($E\approx69$~GPa, $\sigma_y\approx 250$~MPa,
$\rho\approx2700$~kg/m$^3$); carbon-fibre or HY-80 steel changes only the constants.

\subsection{Water-entry survival}\label{sec:entry}
The peak deceleration on water impact is modelled by a monomial fitted from the 2-D interFoam DOE:
\begin{equation}
a_{\text{peak}} = C\, V_{\text{entry}}^a\, d^b\, m_{\text{eff}}^c,
\qquad m_{\text{eff}} = m + m_a, \gpc{monomial fit}
\end{equation}
where $m_a=k_a\rho_w\nabla$ is the axial added mass
($k_a\approx0.03$~\cite{newman2018marine}; a $\sim$3\% correction). The analytic prior from
slamming theory is $a\approx2$ (force $\propto\frac12\rho V^2$ area), $b\approx2$ (diameter$^2$, from area $\propto d^2$),
$c\approx-1$ (Newton). All three exponents are \emph{fitted from CFD} — the prior serves only as a
sanity check. The hard constraint is $a_{\text{peak}}\le 20\,g$; entry angle enters as a secondary
parameter \emph{treated conservatively} (near-vertical worst case) unless the DOE returns strong
angular dependence. Song~et~al.~\cite{song2020waterentry} confirm the $V^2$ scaling for oblique
entry of high-speed projectiles. The GP constraint is monomial: $C V_{\text{entry}}^a d^b
m_{\text{eff}}^c \le a_{\max}$.

\subsection{Wing skew dynamics (residual roll — trajectory gate)}\label{sec:roll}
During the skew from $\Lambda=0^\circ$ to $\Lambda=90^\circ$, an asymmetric wake induces a roll
moment quantified by $C_l(\Lambda)$. This is the single highest-value CFD unknown — currently
unmeasured across a factor of $\sim20\times$ range. The roll-transient time constant is
$\tau\approx0.10$--$0.12$~s. The GP does \emph{not} optimize the skew trajectory
(that is Layer~B / Bayesian optimization); it only imposes the roll-authority feasibility gate:
\begin{equation}
|C_l(\Lambda)| \le C_{l,\max},
\qquad C_{l,\max} \triangleq \frac{\text{tail-fin authority}}{\tfrac12\rho V^2 S b}.
\gpc{monomial bound}
\end{equation}
$C_l(\Lambda)$ is supplied as a fitted posynomial vs $\Lambda$ from the AERO CFD campaign, and
$C_{l,\max}$ is computed from fin geometry: the cruciform ``VFin'' + ``HFin'' pair
($\sim$0.20~m span, $\sim$0.05~m chord) provide a roll-authority moment arm through a
lifting-line fin model. Feedforward tail-deflection scheduling further relaxes the gate (pre-program a
cancelling deflection per the known $C_l(\Lambda)$ function).

\subsection{Static stability and fin sizing}\label{sec:stability}
A trans-medium vehicle without a static-stability constraint could converge to an unflyable optimum.
Renilson~\cite{renilson2018submarine,fossen2011handbook} (Ch.~3, Eq.~113, Table~3.8) defines dimensionless stability
indices in the vertical ($G_V$) and horizontal ($G_H$) planes:
\begin{align}
G_V &= 1 - \frac{M'_w\,(m' + Z'_q)}{M'_q\,Z'_w},
\qquad
G_H = 1 - \frac{N'_v\,(m' - Y'_v)}{N'_r\,Y'_v}\ \text{(by analogy)},\ \gpc{signomial (mixed derivative signs); $G_H$ inferred from vertical-plane form, not in Renilson OCR}\\[4pt]
\text{acceptable:}&\quad G_V\in[0.5,\,0.8],\quad G_H\in[0.2,\,0.4] \qquad\text{(Ray~et~al.~2008)}.
\end{align}
Primes denote non-dimensional stability derivatives divided by $\frac12\rho V L^2$ (e.g.\
$Z'_w = Z_w/(\frac12\rho V L^2)$). The derivatives themselves ($Z'_w, M'_w, Z'_q, M'_q$ for the
vertical plane; $Y'_v, N'_v, N'_r$ for the horizontal plane) can be approximated from fin geometry
via low-aspect-ratio lifting-surface theory~\cite{hoerner1985fluid}: for a cruciform fin set at the
tail, $Z'_w \propto (S_{\text{fin}}/S)(l_{\text{fin}}/\bar{c})$ — the tail volume
coefficient~\cite{hoerner1985fluid} (Ch.~XI),
\begin{equation}
V_H = \frac{S_{\text{fin}}}{S}\,\frac{l_{\text{fin}}}{\bar{c}},
\qquad
C_{L,\text{fin}} \propto V_H\,\alpha_{\text{fin}}. \gpc{monomial}
\end{equation}
For the GP, the $G_V$ and $G_H$ constraints enter as \emph{feasibility gates} with the derivatives
expressed as monomials in fin area and moment arm, using the analytic low-AR lift slope
$\partial C_L/\partial\alpha \approx \pi A_{\text{fin}}/2$~\cite{hoerner1985fluid} ($A_{\text{fin}}
\approx 2\cdot\text{span}_f/\text{chord}_f$). This is the single most important block missing from
the original sizing model. At our small scale ($L\approx1$~m), the cruciform fin pair
(span $\sim\!0.20$~m, chord $\sim\!0.05$~m) gives $A_{\text{fin}}\approx4$, yielding adequate
control authority; the stability gate validates this rather than leaving it unchecked.

\subsection{Propeller efficiency decomposition}\label{sec:propeller}
The GP currently uses a lumped propeller efficiency $\eta_{\text{prop}}\approx0.81$
\cite{allen2000propulsion}. Renilson~\cite{renilson2018submarine} (Ch.~5, Eq.~5.4) decomposes the
quasi-propulsive coefficient (QPC) into physically distinct factors:
\begin{equation}
\text{QPC} = \eta_H \cdot \eta_O \cdot \eta_R,
\qquad \eta_H = \frac{1-t}{1-w}, \gpc{monomial}
\end{equation}
where $t$ = thrust deduction fraction, $w$ = Taylor wake fraction (hull-form dependent,
$\eta_H\approx1.05$ for streamlined bodies~\cite{renilson2018submarine}), $\eta_O$ = open-water
propeller efficiency (function of advance ratio $J$ and loading $K_T/J^2$), and $\eta_R$ = relative
rotative efficiency $\approx1.05$ for single-propeller configurations
($D_{\text{prop}}/D_{\text{hull}}\approx0.4$--$0.7$). The open-water propeller coefficients are
canonical monomials~\cite{newman2018marine} (Ch.~2):
\begin{equation}
K_T = \frac{T}{\rho n^2 D^4},\quad
K_Q = \frac{Q}{\rho n^2 D^5},\quad
J   = \frac{U_A}{n D},\quad
\eta_O = \frac{J}{2\pi}\frac{K_T}{K_Q}. \gpc{monomial}
\end{equation}
If propeller diameter $D$ and RPM $n$ become GP variables, $K_T$ and $K_Q$ are monomials that link
thrust/torque to $D$ and $n$; the design-space constraint $K_T/J^2 \le (K_T/J^2)_{\max}$ prevents
overloading. For the fixed-geometry (``hybrid'') baseline, treat $\eta_O$, $\eta_H$, $\eta_R$ as
constants and revert to $\text{QPC}\approx0.81$. Common AUV system-level values:
motor $\eta_m\approx0.85$ (brushless DC), QPC $\approx0.8$--$0.9$~\cite{renilson2018submarine},
yielding the chain $\eta_{\text{prop}} = \eta_m\eta_{\text{QPC}} \approx 0.7$--$0.75$.

\subsection{Summary of the GP}
The full GP (Layer~A) is:
\begin{equation}\label{eq:gpsummary}
\begin{array}{ll}
\text{minimize}    & m_{\text{pay}}^{-1} \\
\text{subject to} & \text{Eqs. (steady flight, drag, energy, mass, buoyancy/$\xi$, hydro drag,} \\
                  & \hspace{2pt}\text{cavitation, hoop stress, buckling [fitted], entry, roll gate,} \\
	                  & \hspace{2pt}\text{stability $G_V,G_H$, propeller QPC)} \\
\text{with}        & \text{$\xi$ swept; $\Lambda=\text{fixed}$ (solved at each $\Lambda$)} \\
                  & C_{Dp}(\Lambda),\ k(\Lambda),\ C_{L,\text{uw}},\ k_{\text{uw}},\ -C_{p,\min},\ C_l(\Lambda),\ a_{\text{peak}} \text{ from CFD}.
\end{array}
\end{equation}

\section{Model Justification}\label{sec:justify}

\subsection{Why a geometric program, and why sweep the buoyancy state}
The choice of a geometric program is not a matter of convenience but of trust. Because the log
transform makes the problem convex, the solver certifies global optimality and reports, at no extra
cost, how hard each constraint pushes on the objective; a gradient-based optimiser on the same
non-convex problem would offer neither guarantee. The price---writing every relation as a
posynomial---is affordable here because the governing physics is already close to power-law form:
drag scales as a power of Reynolds number, buoyancy as a product of geometric dimensions, and induced
drag as the square of lift.

The buoyancy imbalance $\xi=(W-B)/B$ measures how far the vehicle's weight departs from the weight of
water it displaces. At $\xi=0$ the vehicle is neutrally buoyant and hovers with no trim force; when
$\xi>0$ it is heavy and the stowed wing must generate upward trim lift to hold depth; when $\xi<0$ it
is light and the wing trims downward. The imbalance is swept as a parameter rather than left free for
a specific reason: with $\xi$ fixed, the balance $mg=B(1+\xi)$ is monomial and the whole model
remains a pure GP; freeing $\xi$ would make that equality signomial and would require the slower
sequential (signomial-programming) solver, a route applied across aircraft design~\cite{kirschen2018application}. Sweeping $\xi$ is also the physically informative choice,
because it traces the payload-versus-buoyancy frontier directly: a heavier vehicle carries a larger
mass budget and therefore more payload, up to the point where the wing can no longer lift it in air.

\subsection{Analytic drag with CFD calibration, not a fixed CFD polar}
It would be tempting to insert the measured deployed-wing drag polar directly into the GP. That was
rejected for two reasons. First, a fixed polar is tied to one geometry; once the hull and wing become
design variables (Sec.~\ref{sec:findings}), the drag must scale with them, which only the analytic
build-up---skin friction times a form factor times a wetted-area ratio---can do. Second, and more
subtly, the deployed-wing efficiency inferred from the sweep mesh is not converged: the fitted
lift-curve slope of $0.835$~rad$^{-1}$ is about five times below the thin-airfoil expectation for the
wing's aspect ratio, and the fitted induced factor implies an Oswald efficiency of only $0.07$
against an expected $0.8$--$0.9$. Both discrepancies point to the same cause---a $57{,}000$-cell
sweep mesh under-resolving the wing's response to angle of attack---so the absolute efficiency is
better taken from the analytic prior $e\approx0.85$ than from the mesh-limited fit. What the CFD
\emph{does} deliver reliably are the dimensionless \emph{scalings} of
Eqs.~\eqref{eq:clalpha}--\eqref{eq:kschedule}, the stowed polar (whose aspect-ratio collapse is a
real effect, not a mesh artefact), the cavitation gate, and the cross-validated form factor; these
are what the model uses.

\subsection{Fixed and free geometry, the feasibility gates, and the exclusion of entry}
The model is solved in two configurations. In the first the geometry is fixed at the as-built CAD
dimensions, which verifies that the existing vehicle can meet the mission; in the second the hull
diameter and length, the wing planform, and the fin area are freed within a launch envelope, which
reveals the design the optimiser would prefer. The static-stability requirement, which in its exact
form involves stability derivatives of mixed sign and is therefore signomial, is imposed through the
standard tail-volume coefficients $V_H$ and $V_V$; requiring these to exceed conventional minima is a
GP-legal proxy that guarantees the same static-margin behaviour without leaving the convex world.
Finally, water-entry survival is kept out of the optimisation for two reasons: the peak-deceleration
coefficient is still pending its own two-dimensional impact campaign, and the impact constraint is a
pass/fail survival check that does not size the vehicle in the way the mission-energy and buoyancy
constraints do. It is therefore evaluated after the solve, as an audit of the sized design, rather
than as a driver within it.

\section{Findings}\label{sec:findings}

\subsection{CFD-fitted coefficients}
The fitted coefficients are physically self-consistent and, where an analytic prior exists, they
recover it. Figure~\ref{fig:polars} shows the drag polars for all ten skew angles: the induced-drag
factor grows smoothly and monotonically as the wing folds, with no anomalous point, and every
per-wing polar is reproduced with a coefficient of determination above $0.98$.
Figure~\ref{fig:clalpha} plots the lift-curve slope against skew and overlays the monomial fit of
Eq.~\eqref{eq:clalpha}; the near-unity exponent confirms the $\cos\Lambda$ law.
Figure~\ref{fig:kcollapse} shows the aspect-ratio collapse of the induced factor on a logarithmic
axis, and Fig.~\ref{fig:cavitation} shows why the cavitation data separate into two branches---the
stowed configuration sits at near-zero lift, while the deployed configuration spans a real lift
range---so that fitting them as a single curve, as a first attempt did, is incorrect; the stowed
branch supplies the gate value.

\begin{figure}[h]\centering
\includegraphics[width=0.86\linewidth]{aero_polars_all.png}
\caption{Aerial drag polars for all ten wing-skew angles. The induced-drag factor
$k$ (legend) rises monotonically from $0.64$ at full deployment ($\Lambda=0^\circ$) to $22.9$ when
stowed ($\Lambda=90^\circ$).}
\label{fig:polars}
\end{figure}

\begin{figure}[h]\centering
\includegraphics[width=0.72\linewidth]{sched_CLalpha.png}
\caption{Lift-curve slope versus wing skew. The fitted exponent of $1.105$ confirms the
$C_{L\alpha}\propto\cos\Lambda$ sweep-independence prior to within eleven percent.}
\label{fig:clalpha}
\end{figure}

\begin{figure}[h]\centering
\includegraphics[width=0.72\linewidth]{sched_k.png}
\caption{Induced-drag factor versus wing skew on a logarithmic axis: the thirty-six-fold rise is the
signature of the aspect-ratio collapse as the wing folds flat.}
\label{fig:kcollapse}
\end{figure}

\begin{figure}[h]\centering
\includegraphics[width=0.80\linewidth]{cavitation_split.png}
\caption{Cavitation suction peak against lift coefficient. The stowed configuration (near zero lift)
and the deployed configuration (finite lift range) form two distinct branches; the stowed branch,
$-C_{p,\min}\approx1.16$, sets the cavitation gate.}
\label{fig:cavitation}
\end{figure}

\subsection{The payload--buoyancy frontier}
Figure~\ref{fig:payload} traces payload against the buoyancy imbalance for both geometry cases. In
both, payload rises with $\xi$, because a heavier vehicle carries a larger mass budget; in both, the
frontier terminates where the aerial lift coefficient reaches its stall ceiling and the wing can no
longer support the vehicle in air. The fixed-geometry vehicle carries between $2.4$ and $3.5$~kg of
payload across the feasible range of $\xi$, which confirms that the as-built design meets a several-
kilogram payload goal. Freeing the geometry roughly doubles the payload, to between $5.9$ and
$8.4$~kg, for a reason that Sec.~\ref{sec:optgeom} makes precise: the optimiser enlarges the hull to
the edge of the launch envelope and thereby buys buoyancy.

\begin{figure}[h]\centering
\includegraphics[width=0.72\linewidth]{payload_xi.png}
\caption{Payload versus buoyancy imbalance $\xi$ for the fixed CAD geometry and the
envelope-constrained free geometry. Both curves end where the aerial lift coefficient reaches the
$0.60$ stall ceiling.}
\label{fig:payload}
\end{figure}

\subsection{The optimised geometry and what constrains it}\label{sec:optgeom}
Table~\ref{tab:soln} lists the free-geometry solution across the imbalance sweep. The pattern is
uniform and revealing: at every value of $\xi$ the hull diameter and length pin to their
envelope caps of $0.12$~m and $1.30$~m, the wing span pins to the geometric limit of one hull length
minus $0.30$~m, and the aerial lift coefficient pins to its $0.60$ stall ceiling. In other words the
optimiser drives the vehicle to three limits at once---how large the launch platform allows it to be,
how wide the fuselage allows the wing to span, and how hard the wing can work in air---and the
payload it returns is whatever remains after the mission energy and the pressure hull have been paid
for. The cruciform fins, sized by the static-stability gate, come out at about $86$~cm$^2$ of total
area on a $0.585$~m moment arm, adding roughly three-quarters of a kilogram; the roll-authority gate
is satisfied with margin by the fins that stability already demands. Neither gate drives the payload,
which is exactly the behaviour a feasibility gate should show.

\begin{table}[h]\centering
\caption{Free-geometry Layer-A solution across the buoyancy sweep. The row at $\xi=0.15$ is the
reference design point.}\label{tab:soln}
\small
\begin{tabular}{@{}rrrrrrrr@{}}\toprule
$\xi$ & payload & mass & $d$ & $L$ & span & $S$ & $S_{\text{fin}}$ \\
      & (kg) & (kg) & (mm) & (mm) & (m) & (m$^2$) & (cm$^2$) \\\midrule
$0.00$ & $5.86$ & $9.12$ & $120$ & $1300$ & $1.00$ & $0.097$ & $64.8$ \\
$0.05$ & $6.28$ & $9.57$ & $120$ & $1300$ & $1.00$ & $0.102$ & $71.4$ \\
$0.10$ & $6.71$ & $10.03$ & $120$ & $1300$ & $1.00$ & $0.107$ & $78.4$ \\
$0.15$ & $7.14$ & $10.48$ & $120$ & $1300$ & $1.00$ & $0.112$ & $85.7$ \\
$0.20$ & $7.56$ & $10.94$ & $120$ & $1300$ & $1.00$ & $0.117$ & $93.3$ \\
$0.25$ & $7.99$ & $11.39$ & $120$ & $1300$ & $1.00$ & $0.122$ & $101.2$ \\
$0.30$ & $8.41$ & $11.85$ & $120$ & $1300$ & $1.00$ & $0.127$ & $109.5$ \\\bottomrule
\end{tabular}
\end{table}

\subsection{Sensitivities and design levers}
The dual solution ranks the levers that matter. The strongest is the displaced volume, with a
sensitivity of $-1.43$: enlarging the hull is the most direct way to increase payload, which is why
the optimiser pins the hull to its bounds. Battery specific energy follows at $-0.49$, then aerial
efficiency and range at about $\pm0.30$, and required underwater speed at $+0.38$. The stowed
induced-drag factor registers a sensitivity of only $+0.001$---negligible---because the stowed wing
operates at almost zero lift and its trim drag is correspondingly tiny. This last number is worth
dwelling on, because it is the quantitative form of the vehicle's design thesis: the folded wing
costs almost nothing in drag, so the penalty for keeping the wing rather than discarding it before
water entry is small.

\subsection{Verification of the surrogates and the sized model}
Each surrogate and the assembled program were checked by five independent means before either was
trusted. The first is the air--water cross-validation already noted: the stowed induced-drag factor
agrees to seven percent between the two campaigns, a comparison that no single campaign could make.
The second is agreement with analytic priors: the lift-slope exponent of $1.105$ against the value of
unity in Eq.~\eqref{eq:clalpha}, the Reynolds exponent of $-0.112$ bracketed by the skin-friction and
form-drag limits, and the near-zero exponent of the cavitation peak all fall where first principles
say they should. The third is reconstruction quality: every per-wing polar is recovered with a
coefficient of determination above $0.98$ and every retained monomial fit with a root-mean-square
logarithmic error below $0.02$. The fourth is GP-legality, verified term by term: each quantity that
enters a constraint has a positive-coefficient form, the one exception---the cambered deployed wing,
whose zero-lift intercept is negative---being replaced by a positive softmax-affine surrogate. The
fifth is physical monotonicity: lift rises with angle of attack, drag rises as the wing deploys, and
the suction peak is speed-independent, at every point sampled. The assembled program was audited in
the same spirit: its dual sensitivities are all of the expected sign and magnitude, and its binding
constraints---mass closure and buoyancy---are exactly those a near-neutral trans-medium vehicle
should be limited by, which is itself evidence that the model is wired correctly.

\section{Layer-A Conclusions}

Three conclusions follow from the Layer-A study. First, the vehicle is buoyancy- and
envelope-limited rather than aerodynamically or structurally limited: the binding constraints are the
hull-size caps and the aerial stall ceiling, and the dominant sensitivity is to displaced volume, so
payload is bought almost directly with the size the launch platform will allow. The as-built CAD
vehicle carries roughly three kilograms of payload; an envelope-optimised vehicle of the same family
carries about seven. Second, the design thesis survives contact with the numbers: the stowed wing's
trim drag is negligible in the sizing, so retaining a steerable wing through water entry, rather than
jettisoning it, imposes almost no payload penalty, and the lift-curve slope's clean $\cos\Lambda$
scaling confirms that the oblique wing behaves as classical sweep theory predicts. Third, the model
is honest about where it is weakest: the absolute efficiency of the deployed wing is not yet
grid-converged and is therefore carried through an analytic prior rather than the mesh-limited fit,
the drag polar over the measured lift range is too narrow to support a full softmax-affine surrogate,
and the water-entry constant awaits its own campaign. These are the three inputs a refined study
should target---a grid-convergence mesh for the deployed wing, a wider angle-of-attack sweep or a
section-level profile-drag polar, and the two-dimensional impact campaign---after which the same
Layer-A program, unchanged in structure, will return sharper numbers. The value of the present model
is that it already tells the designer where to look: at the size of the launch envelope, the energy
density of the battery, and the efficiency of the deployed wing, in that order.

% \bibliographystyle is set by new-aiaa.cls
\bibliography{uauv}

\end{document}
